
\medskip
\phantomsection\label{sga3:exp5:lemma:1.2}
\noindent\textbf{Lemma 1.2.}---
\textit{Let \(d_0,d_1:X_1\rightrightarrows X_0\) be a diagram in
\((\mathbf{Sch})\), and let \((Y,p)\) be its cokernel in
\((\mathbf{Esp.An})\).
\begin{enumerate}
\item[(i)]\phantomsection\label{lem:sga3-V-1.2-i}%
If \(Y\) is a scheme and \(p\) is a morphism of schemes,
  then \((Y,p)\) is a cokernel of \((d_0,d_1)\) in
  \((\mathbf{Sch})\).
\item[(ii)]\phantomsection\label{lem:sga3-V-1.2-ii}%
Suppose that every point of \(X_0\) has a saturated open
  neighborhood \(U\) with the following property. Let
  \(\widetilde d_0\) and \(\widetilde d_1\) be the restrictions of
  \(d_0\) and \(d_1\) to
  \(d_0^{-1}(U)=d_1^{-1}(U)\), and let \((Q,q)\) be the cokernel of
  \((\widetilde d_0,\widetilde d_1)\) in \((\mathbf{Esp.An})\). If
  \(Q\) is a scheme and \(q\) is a morphism of schemes, then
  \((Y,p)\) is a cokernel of \((d_0,d_1)\) in
  \((\mathbf{Sch})\).
\end{enumerate}}

\hyperref[lem:sga3-V-1.2-i]{Part~(i)} is proved in \hyperref[proof:sga3-V-4.1-c]{\S4(c)}; since the proof is short, let us repeat
it here. Let \(f:X_0\to Z\) be a morphism of schemes such that
\(fd_0=fd_1\). By hypothesis there is a unique morphism of ringed
spaces \(r:Y\to Z\) such that \(f=rp\). It remains to show that, for
every \(y\in Y\), the homomorphism
\(\mathscr O_{r(y)}\to\mathscr O_y\) induced by \(r\) is local. This
follows because \(p\) is surjective, so \(y=p(x)\) for some \(x\), and
because the homomorphism
\(\mathscr O_{f(x)}\to\mathscr O_x\) induced by \(f\) is local.

\hyperref[lem:sga3-V-1.2-ii]{Part~(ii)} follows from \hyperref[lem:sga3-V-1.2-i]{part~(i)} and \hyperref[lem:sga3-V-1.1]{the preceding lemma}.

\medskip
\phantomsection\label{par:sga3-V-1-c}%
\noindent\textbf{c)} In this expos\'e we study the existence of
\(\operatorname{Coker}(d_0,d_1)\) when the double arrow
\((d_0,d_1)\) is placed in a richer context. More precisely, let
\[
  X_2=X_1\times_{d_1,d_0}X_1
\]
be the fiber product of the diagram
\begin{equation}\SGAInlineEquationTarget{sga3:v:lemma:1-2:equation:star}
  \resizebox{0.20\textwidth}{!}{%
    \SGAThreeNativeInput{figures/exp5/p253_star_fiber_product.png}}
  \tag{*}\label{fig:sga3-exp5-star-fiber-product}
\end{equation}
and let \(d'_0,d'_2:X_2\to X_1\) be the two canonical projections.
Thus, by definition, the square
\begin{equation}\SGAInlineEquationTarget{sga3:v:lemma:1-2:equation:0}
  \resizebox{0.21\textwidth}{!}{%
    \SGAThreeNativeInput{figures/exp5/p253_square_0.png}}
  \tag{0}\label{fig:sga3-exp5-square-0}
\end{equation}
is Cartesian.

In addition, suppose that a third arrow \(d'_1:X_2\to X_1\) is
given. We say that
\[
  (d_0,d_1:X_1\rightrightarrows X_0,d'_1)
\]
is a \(\mathscr C\)-\emph{groupoid} if, for every object \(T\) of
\(\mathscr C\), \(X_1(T)\) is the set of arrows of a groupoid
\(X_\ast(T)\) whose set of objects is \(X_0(T)\), whose source map is
\(d_1(T)\), whose target map is \(d_0(T)\), and whose composition map
is \(d'_1(T)\). As usual, we identify
\[
 (X_1\times_{d_1,d_0}X_1)(T)
 \simeq X_1(T)\times_{d_1(T),d_0(T)}X_1(T),
\]
and recall that a groupoid is a category in which every arrow is
invertible.\footnote{\label{note:sga3-V-5}Editor's note: in this case \(X_2(T)\) is the set
of pairs \((f_2,f_1)\) of composable arrows, that is, pairs such that
\(d_0(f_1)=d_1(f_2)\); the maps \(d'_0,d'_1,d'_2\) send
\((f_2,f_1)\), respectively, to \(f_2,f_2\circ f_1,f_1\).}

If \(\varphi\) is an arrow of the groupoid \(X_\ast(T)\), the map
\(f\mapsto\varphi\circ f\) is a bijection from the set of arrows \(f\)
whose target equals the source of \(\varphi\) onto the set of arrows
having the same target as \(\varphi\). This is equivalently expressed
by saying that the square
\begin{equation}\SGAInlineEquationTarget{sga3:v:lemma:1-2:equation:1}
  \resizebox{0.19\textwidth}{!}{%
    \SGAThreeNativeInput{figures/exp5/p254_square_1.png}}
  \tag{1}\label{fig:sga3-exp5-square-1}
\end{equation}
is Cartesian.

Similarly, \(g\mapsto g\circ\varphi\) is a bijection from the set of
arrows \(g\) of \(X_\ast(T)\) whose source equals the target of
\(\varphi\) onto the set of arrows having the same source as
\(\varphi\). Equivalently, the square
\begin{equation}\SGAInlineEquationTarget{sga3:v:lemma:1-2:equation:2}
  \resizebox{0.19\textwidth}{!}{%
    \SGAThreeNativeInput{figures/exp5/p254_square_2.png}}
  \tag{2}\label{fig:sga3-exp5-square-2}
\end{equation}
is Cartesian.

On the other hand, let \(s:X_0\to X_1\) be the unique arrow of
\(\mathscr C\) such that, for every \(T\), the map
\(s(T):X_0(T)\to X_1(T)\) sends each object of \(X_\ast(T)\) to its
identity arrow.\footnote{\label{note:sga3-V-6}Editor's note: \(T\mapsto s(T)\) defines an
element of \(\operatorname{Hom}(h_{X_0},h_{X_1})\), and the latter is
equal to \(\operatorname{Hom}(X_0,X_1)\) by the Yoneda lemma.}
The arrow \(s\) satisfies
\begin{align}
  d_1s&=\operatorname{id}_{X_0},
    \tag{3}\label{eq:sga3-exp5-3}\\
  d_0s&=\operatorname{id}_{X_0}.
    \tag{3 bis}\label{eq:sga3-exp5-3bis}
\end{align}
Finally, associativity of the composition maps \(d'_1(T)\) is
equivalent to commutativity of the diagram
\begin{equation}\SGAInlineEquationTarget{sga3:v:lemma:1-2:equation:4}
  \resizebox{0.72\textwidth}{!}{%
    \SGAThreeNativeInput{figures/exp5/p254_associativity_4.png}}
  \tag{4}\label{fig:sga3-exp5-associativity-4}
\end{equation}

Conversely, conditions
\hyperref[sga3:v:lemma:1-2:equation:1]{(1)},
\hyperref[sga3:v:lemma:1-2:equation:2]{(2)}, and
\hyperref[sga3:v:lemma:1-2:equation:4]{(4)}, together with the
existence of an arrow \(s\) satisfying
\hyperref[eq:sga3-exp5-3]{(3)}, imply that
\((X_1\rightrightarrows X_0,d'_1)\) is a
\(\mathscr C\)-groupoid. Condition
\hyperref[eq:sga3-exp5-3]{(3)} is harmless: it merely ensures that
\(d_1(T):X_1(T)\to X_0(T)\) is surjective for every
\(T\in\mathscr C\). In the sequel we shall mainly use the Cartesian
squares
\hyperref[sga3:v:lemma:1-2:equation:0]{(0)},
\hyperref[sga3:v:lemma:1-2:equation:1]{(1)}, and
\hyperref[sga3:v:lemma:1-2:equation:2]{(2)}, summarized in the diagram
\begin{equation}\SGAInlineEquationTarget{sga3:v:lemma:1-2:equation:0-1-2}
  \resizebox{0.34\textwidth}{!}{%
    \SGAThreeNativeInput{figures/exp5/p255_diagram_012.png}}
  \tag{0,1,2}\label{fig:sga3-exp5-diagram-012}
\end{equation}
The two left-hand squares in this diagram (that is, squares
\hyperref[sga3:v:lemma:1-2:equation:0]{(0)} and
\hyperref[sga3:v:lemma:1-2:equation:2]{(2)}) are Cartesian; the first row
is exact; and \(X_2\) is identified with the fiber product
\(X_1\times_{d_0,d_0}X_1\).

We use associativity only indirectly, for example to ensure the
existence of an arrow \(s\) satisfying
\hyperref[eq:sga3-exp5-3]{(3)} and
\hyperref[eq:sga3-exp5-3bis]{(3 bis)}, or to ensure the existence of
an arrow
\begin{equation}
  \sigma:X_1\longrightarrow X_1,
  \qquad d_0\sigma=d_1,\quad d_1\sigma=d_0.
  \tag{\ensuremath{\dagger}}\label{eq:sga3-exp5-dagger}
\end{equation}
Here \(\sigma\) is chosen so that
\(\sigma(T):X_1(T)\to X_1(T)\) sends every arrow of \(X_\ast(T)\) to
its inverse.\footnote{\label{note:sga3-V-7}Editor's note: the Yoneda lemma shows that
\(\sigma\) is an involutive automorphism of \(X_1\); this will be used,
for example, in \hyperref[par:sga3-V-3-e]{3(e)} and in \hyperref[thm:sga3-V-4.1]{Theorem~4.1}.}

By an abuse of language, we shall sometimes call a diagram
\begin{equation*}
  \resizebox{0.40\textwidth}{!}{%
    \SGAThreeNativeInput{figures/exp5/p255_groupoid_triple_arrow.png}}
  \phantomsection\label{fig:sga3-exp5-groupoid-triple-arrow}
\end{equation*}
a \(\mathscr C\)-groupoid when
\hyperref[sga3:v:lemma:1-2:equation:0]{(0)},
\hyperref[sga3:v:lemma:1-2:equation:1]{(1)}, and
\hyperref[sga3:v:lemma:1-2:equation:2]{(2)} are Cartesian,
\hyperref[sga3:v:lemma:1-2:equation:4]{(4)} is commutative, and
there exists an \(s\) satisfying
\hyperref[eq:sga3-exp5-3]{(3)}. Thus \(X_2\) may be ``a'' fiber
product of \hyperref[sga3:v:lemma:1-2:equation:star]{(*)} without
being ``the'' fiber product of
\hyperref[sga3:v:lemma:1-2:equation:star]{(*)}.\footnote{\label{note:sga3-V-8}Editor's
note: see \hyperref[par:sga3-V-2-a]{Example~2(a)} below.}

\medskip
\phantomsection\label{sga3:exp5:terminology:c-groupoid}
\noindent\textit{Terminology.} In place of the
\(\mathscr C\)-groupoid \(X_\ast\), we shall also speak of the
\emph{groupoid \(X_\ast\) with base \(X_0\)}, or of the
\emph{pre-equivalence relation \(X_\ast\) in \(X_0\)}.

\section*{2. Examples of \(\mathscr C\)-groupoids}
\phantomsection\label{sec:sga3-V-2}
\addcontentsline{toc}{section}{2. Examples of
  \texorpdfstring{\(\mathscr C\)}{C}-groupoids}

\phantomsection\label{par:sga3-V-2-a}%
\noindent\textbf{a)} Let \(X\) be an object of \(\mathscr C\), and let
\(G\) be a \(\mathscr C\)-group acting on the left on \(X\). We denote
by \(d_0:G\times X\to X\) the arrow defining the action of \(G\) on
\(X\), by \(d_1:G\times X\to X\) the projection onto the second
factor, and by \(\mu:G\times G\to G\) the arrow defining
